\section{The fluid problem behind the envelope equation}
\label{sec:fluid}
\subsection{Why Laplace's equation governs the fluid interior}
Let $\boldsymbol{u}=(u,v,w)$ denote velocity. Incompressibility is
$\nabla\cdot\boldsymbol{u}=0$; irrotationality is
$\nabla\times\boldsymbol{u}=0$. For a smooth flow in a simply connected
domain, the latter permits $\boldsymbol{u}=\nabla\phi$. Therefore
\begin{equation}\label{eq:laplace}
 \nabla^2\phi=0,\qquad z<\zeta(x,y,t).
\end{equation}
This is (2.1). It is a constraint at each time, not a wave equation in the
fluid interior. Time evolution enters through the surface conditions.

\subsection{Kinematic condition: the surface moves with the water}
Write the surface as $F=z-\zeta(x,y,t)=0$. A surface fluid particle must
remain there, so $DF/Dt=F_t+\boldsymbol{u}\cdot\nabla F=0$. Thus
\begin{equation}\label{eq:kinematic}
 \zeta_t+\phi_x\zeta_x+\phi_y\zeta_y=\phi_z,
 \qquad z=\zeta.
\end{equation}
The left side includes horizontal advection of a sloping surface.
$\zeta_t=\phi_z$ alone is the linearized condition.

\subsection{Dynamic condition: atmospheric pressure at the surface}
Integration of the inviscid potential-flow momentum equation gives
\[
 \phi_t+\frac12|\nabla\phi|^2+\frac p\rho+gz=C(t).
\]
An additive time-dependent change to $\phi$ absorbs $C(t)$. Subtract
constant atmospheric pressure and restrict to the surface:
\begin{equation}\label{eq:bernoulli}
 \phi_t+\frac12|\nabla\phi|^2+g\zeta=0,
 \qquad z=\zeta.
\end{equation}
This is Bernoulli's equation. Inviscid does not mean linear:
$|\nabla\phi|^2$ remains.

Dysthe's (2.2) differentiates Bernoulli along a surface fluid particle.
For smooth $f$, $Df/Dt=f_t+\nabla\phi\cdot\nabla f$.
Using the particle's vertical velocity $\phi_z$ gives
\begin{equation}\label{eq:dysthe-dynamic}
 \phi_{tt}+g\phi_z+\partial_t|\nabla\phi|^2
 +\frac12\nabla\phi\cdot\nabla|\nabla\phi|^2=0,
 \qquad z=\zeta.
\end{equation}
Differentiation of $\phi_t$ contributes
$\nabla\phi\cdot\nabla\phi_t$; differentiation of
$\tfrac12|\nabla\phi|^2$ contributes another equal term. Together they
give $\partial_t|\nabla\phi|^2$, explaining its coefficient in (2.2).

\subsection{Why expand the boundary condition about $z=0$?}
For smooth $f$ and small $\zeta$,
\begin{equation}\label{eq:boundary-taylor}
 f(x,y,\zeta,t)=f(x,y,0,t)+\zeta f_z(x,y,0,t)
 +\frac12\zeta^2 f_{zz}(x,y,0,t)+\cdots.
\end{equation}
This moves the equations to a fixed plane while retaining the effects of
surface displacement as extra terms. It does not flatten the actual wave.

In (2.4), $L=\partial_t^2+g\partial_z$. Taylor expansion of
\eqref{eq:dysthe-dynamic} gives
\[
 L\phi+\zeta L\phi_z+\partial_t|\nabla\phi|^2+Q=0,
 \qquad z=0,
\]
where the retained remainder is
\[
 Q=\frac12\zeta^2L\phi_{zz}
 +\frac12\nabla\phi\cdot\nabla|\nabla\phi|^2
 +\zeta\partial_{zt}|\nabla\phi|^2.
\]
These pieces come from another Taylor term, the cubic material derivative,
and the Taylor correction to the quadratic velocity term.

The expanded kinematic condition (2.5) is
\[
 \zeta_t-\phi_z+\nabla_h\cdot(\zeta\nabla_h\phi)
 +\nabla_h\cdot\left(\frac12\zeta^2\nabla_h\phi_z\right)=0,
 \qquad z=0,
\]
at the indicated order, where $\nabla_h=(\partial_x,\partial_y,0)$.
For example, its quadratic terms are initially
$\nabla_h\phi\cdot\nabla_h\zeta-\zeta\phi_{zz}$.
Since $-\phi_{zz}=\nabla_h^2\phi$, their sum becomes
$\nabla_h\cdot(\zeta\nabla_h\phi)$. The divergence form is an algebraic
reorganization, not an extra hypothesis.

\subsection{A first nonlinear calculation: the second Stokes harmonic}
A \emph{Stokes wavetrain} is a periodic travelling solution of the nonlinear
water-wave equations. Its surface is not exactly sinusoidal. Here is a short
calculation showing how a bound harmonic enters.

Suppress $y$ and let $\theta=k_0x-\omega_0t$, with
$\omega_0^2=gk_0$. The leading real potential is
\[
 \phi_1=\frac{a\omega_0}{k_0}e^{k_0z}\sin\theta.
\]
Through second order, try a bound elevation harmonic
$\zeta=a\cos\theta+C\cos(2\theta)$ and no second potential harmonic.
Expand Bernoulli at $z=0$. Its second-order part is
\[
 -ga^2k_0\cos^2\theta+\frac12a^2\omega_0^2
                         +gC\cos(2\theta)=0.
\]
Since $\cos^2\theta=(1+\cos2\theta)/2$, this gives
\begin{equation}\label{eq:stokes-second}
 C=\frac12k_0a^2,\qquad
 \zeta=a\cos\theta+\frac12k_0a^2\cos(2\theta)+\cdots.
\end{equation}
The kinematic condition independently gives
$2\omega_0C=k_0\omega_0a^2$, the same result. At a crest both elevation
terms are positive, while at a trough the second term is still positive;
the profile has sharper crests and flatter troughs. The harmonic oscillates
at $2\omega_0$, whereas a free linear mode of wavenumber $2k_0$ has
frequency $\sqrt2\omega_0$. It is therefore forced by the carrier,
not an independent linear wave. This simple calculation illustrates
the Stokes comparison mentioned immediately after (2.13).

\paperlink{Read (2.1)--(2.5), printed pp.~106--107.
Connect each equation to incompressibility, surface motion, or pressure.}
\checkpoint{Derive \eqref{eq:kinematic} from $DF/Dt=0$ and identify
which terms disappear upon linearization.}
\clearpage
